\subsection{Stationary VAR Generation: Heaps' Parameterization}
\label{sec:implemented-heaps}

\subsubsection{Desired Property and Mathematical Conditions}
We want linear multivariate series whose statistical regime is the same
throughout generation, including the first returned generated value. For a
fixed Gaussian VAR, this means choosing stable coefficients and drawing
the whole initial lag history from its stationary joint distribution.
In practical terms, the mean, variance, and correlations at a fixed time
separation remain unchanged when the observation window moves. Individual
trajectories still fluctuate because each step receives a new innovation.

\paragraph{Why the method is called Heaps.}
The name credits Sarah E. Heaps' symmetric parameterization and reconstruction
procedure, which build on the earlier mapping of
\citet{ansley1986reparameterizing}. Here, \emph{symmetric} describes the matrix
square roots used in that procedure. The reconstructed coefficient matrices
can be asymmetric: the influence of series $j$ on series $k$ can differ from
the influence in the opposite direction. Our generator uses this construction
with a project-defined sampling distribution.

\paragraph{Published result used.}
\citet[Sections~2.1--2.2 and Appendix]{heaps2023stationarity} parameterizes
stationary Gaussian VAR($L$) models through intermediate parameters with
simple constraints. A \emph{partial autocorrelation} describes the remaining
linear association between two observations after accounting for the
observations between them. For several series, these associations are
represented by matrices $\mathbf P_\ell^{\mathrm{PAC}}$ built from standardized
prediction errors. Every singular value of each matrix must be below one.
A singular value measures how strongly a matrix scales a direction in its
input space; this condition therefore limits all directions together.

Equation~\eqref{eq:heaps-pac-map} obtains these matrices from freely chosen
$\mathbf U_\ell$ by mapping each singular value $u\geq0$ to
$u/\sqrt{1+u^2}<1$. Given these admissible partial autocorrelations and a
positive-definite innovation covariance $\Sigma_\varepsilon$, the paper's
reverse covariance and coefficient recursions recover stationary VAR
coefficients. Positive definiteness means that every nonzero linear
combination of the innovation components has positive variance.

The resulting causal VAR satisfies $\rho(\mathbf B)<1$: every eigenvalue of
the companion matrix $\mathbf B$ has magnitude below one. This matrix
advances the entire lag history by one step. Its powers tend to zero, so
the effect of changing the initial history disappears when future
innovations are held equal. Together with independent Gaussian innovations,
this gives a unique causal stationary law. Drawing the initial history from
that law also makes the finite generated trajectory stationary from its start.

\subsubsection{Implementation and Stationary Initialization}
\path{src/tsfmxai/generators/heaps.py} implements the following sequence.

\paragraph{Step 1: sample the free parameters.}
The mechanism sampler samples the order $L$, one square
matrix $\mathbf U_\ell$ per lag, and the stationary mean. The recorded
comparison uses orders uniformly drawn from 1--12, independent normal
matrix entries with standard deviation 0.45, and a mean fixed at zero.
The sampler fixes $\Sigma_\varepsilon=0.05^2 I$ from the configured innovation
scale. The lower-level constructor also accepts a user-specified
positive-definite covariance, including correlations between series.
These sampling choices determine which mechanisms occur most often.

\paragraph{Step 2: reconstruct jointly stable coefficients.}
The implementation applies the singular-value map using a singular-value
decomposition. The coefficient reconstruction follows the Appendix's two
recursions: it first recovers prediction-error covariances, then reconstructs
the coefficient matrices $\mathbf C_\ell$ and lagged covariances. This is
necessary because the combined effect of all coefficients determines
stability. The intermediate partial autocorrelations provide coordinates
in which that joint restriction is already satisfied. For one series and
one lag, the result is an AR(1) coefficient satisfying $|c_{1,1,1}|<1$;
the sign of the sampled scalar parameter is preserved.

\paragraph{Step 3: construct the equations and check their covariance.}
The final coefficients are converted to one linear AST per target, including
the intercept required by the selected stationary mean. The constructor
forms the joint lag-state covariance $\mathbf V$ from the reconstructed
autocovariances and checks
\begin{equation}
 \mathbf V=\mathbf B\mathbf V\mathbf B^\top+\mathbf Q,
 \label{eq:implemented-heaps-lyapunov}
\end{equation}
where $\mathbf Q$ inserts $\Sigma_\varepsilon$ into the newest lag block.
The equation balances the covariance propagated from the previous state
with the covariance introduced by the next innovation. For a stable
$\mathbf B$, its solution is
$\mathbf V=\sum_{v=0}^{\infty}\mathbf B^v\mathbf Q(\mathbf B^\top)^v$.
Independent numerical checks require a companion spectral radius below one,
partial-autocorrelation singular values below one, positive-definite
$\mathbf V$, and a relative Lyapunov residual below $10^{-8}$. The residual
is the Frobenius norm of the difference between the two sides of
Equation~\eqref{eq:implemented-heaps-lyapunov}, divided by the Frobenius norm
of $\mathbf V$. It checks reconstruction accuracy in floating-point
arithmetic. A mechanism failing these checks raises an error.

\paragraph{Step 4: draw a compatible history and simulate.}
The entire initial lag history is drawn jointly from the corresponding
Gaussian stationary distribution, using a Cholesky factor of $\mathbf V$.
This factor transforms independent standard-normal draws into correlated
values with the required covariance. Sampling all initial values together
preserves both cross-series and cross-lag correlations. The recurrence then
runs immediately with fresh Gaussian innovations independent across time
and independent of the initial history. Each innovation vector has covariance
$\Sigma_\varepsilon$. No preliminary simulation is needed for this
initialization; stationarity holds in exact arithmetic, subject to the
reported numerical checks. This treatment of initial observations follows
\citet[Section~4 and Appendix, pp.~8 and 11]{heaps2023stationarity}.

\paragraph{Scope of the implementation and evaluation.}
The implementation supports several series and lags; VAR(1) is the single-lag
case. The main trade-off experiment uses one series, while the unit tests
also check multivariate covariance reconstruction, agreement between the
AST and the matrix equation, and joint stationary initialization. The final
equations expose signed linear coefficients and variable--lag support.
Nonlinear products belong to the preceding AST families. Gaussian innovations
have unbounded support, so stationary trajectories can still cross a finite
amplitude threshold. The explosion diagnostic measures these realized
events, and the low-amplitude diagnostic separately measures paths that
remain small within the observed window.
